\begin{proof}[Proof of \cref{obj:lem:separated-value-mixtures}]
\begin{enumerate}
\item \textit{Measurable decision laws.}
Denote the two priors by
\[
\pi_0:=\text{the first prior},
\qquad
\pi_1:=\text{the second prior}.
\]
Write
\[
\mathcal O_d:=[d]\times\{0,1\}^2,
\qquad
\mathcal Z:=\prod_{i=1}^n\mathcal Z_i,
\qquad
\mathcal R:=\text{the public seed space}.
\]
A decision triple has the form
\[
\mathsf w_{\mathrm{dec}}
=(\mathbf z,\mathsf r,\mathsf u)
\in\mathcal Z\times\mathcal R\times\mathbb R.
\]
For an observed record vector \(\mathbf o\in\mathcal O_d^n\) and a seed value \(\mathsf r\in\mathcal R\), define the probability transcript law by iterating the protocol kernels:
\[
\mathsf L_{\mathrm{tr}}(\mathbf o,\mathsf r)(d\mathbf z)
:=
\prod_{i=1}^n
Q_i(dz_i\mid o_i,\mathbf z_{<i},\mathsf r).
\]
Here the product denotes successive kernel integration in participant order. Define a probability measure for each fixed record vector by
\[
\mathsf D_{\mathbf o}(E)
:=
\int_{\mathcal R}\int_{[0,1]}\int_{\mathcal Z}
\mathbf 1_E(\mathbf z,\mathsf r,\mathsf u)\,
\mathsf L_{\mathrm{tr}}(\mathbf o,\mathsf r)(d\mathbf z)\,
d\mathsf u\,\lambda(d\mathsf r),
\]
for every Borel decision event \(E\). Each measure is a probability measure because all the message kernels and the seed law are probability measures.

For every nonnegative measure \(P\) on the full-record space, let \(P_O\) be its observed-record marginal and define
\[
\mathsf L_{\mathrm{dec}}(P)
:=
\sum_{\mathbf o\in\mathcal O_d^n}
\left(\prod_{i=1}^n P_O(\{o_i\})\right)\mathsf D_{\mathbf o}.
\]
This map is measurable: each coefficient is a finite product of measurable evaluations of the observed marginal, and the measures \(\mathsf D_{\mathbf o}\) are fixed.

The auxiliary target on all nonnegative measures uses the same welfare formula, with the following conventions:
\[
P_{\mathrm{real}}(E):=
\begin{cases}
P(E),&P(E)<\infty,\\
0,&P(E)=\infty,
\end{cases}
\qquad
\mu_{aj}(P):=
\frac{P_{\mathrm{real}}\{X=j,Y^a=1\}}
     {P_{\mathrm{real}}\{X=j\}},
\qquad
V(P):=
\frac1d\sum_{j=1}^d\max\{\mu_{0j}(P),\mu_{1j}(P)\}.
\]
Division by zero, including \(1/d\) when \(d=0\), is assigned value zero; empty sums and products have values zero and one, respectively. These formulas make the target measurable and agree with the stated conditional means and welfare on \(\mathcal V_d\).

For a probability law \(P\), the finite atomic decomposition is
\[
P_O^{\otimes n}
=
\sum_{\mathbf o\in\mathcal O_d^n}
\left(\prod_{i=1}^nP_O(\{o_i\})\right)\delta_{\mathbf o}.
\]
Consequently, integrating the protocol against the canonical product sampling law gives exactly \(\mathsf L_{\mathrm{dec}}(P)\). Moreover, \cref{obj:ass:iid-people,obj:ass:independent-randomness} give, for every \(P\in\mathcal V_d\),
\[
\mathcal L_P(\mathbf O,R,U)
=
P_O^{\otimes n}\otimes\lambda
\otimes\operatorname{Uniform}[0,1],
\qquad
\mathsf L_{\mathrm{dec}}(P)
=
\mathcal L_{P,Q}(\mathbf Z,R,U).
\]
Define the completed mixtures, for \(\ell\in\{0,1\}\), by
\[
\mathsf L_\ell(E)
:=
\int \mathsf L_{\mathrm{dec}}(P)(E)\,\pi_\ell(dP).
\]
Causal support and the preceding equality imply
\[
\mathsf L_\ell(E)
=
\int_{\mathcal V_d}
\Pr_{P,Q}\{(\mathbf Z,R,U)\in E\}\,\pi_\ell(dP).
\]
Thus these are the decision mixtures in the statement, and their seed-transcript marginals are \(N_\ell\).
% lean: CausalSmith.Stat.LdpOptvalueUniformFrontier.separated_value_mixtures

\item \textit{Independence of the analyst randomizer.}
For \(P\in\mathcal V_d\), the preceding construction gives
\[
\mathsf L_{\mathrm{dec}}(P)
(\mathcal Z\times\mathcal R\times\mathbb R)=1.
\]
For Borel sets \(E_Z\subseteq\mathcal Z\), \(E_R\subseteq\mathcal R\), and \(E_U\subseteq\mathbb R\), successive integration yields
\[
\begin{aligned}
&\mathsf L_{\mathrm{dec}}(P)(E_Z\times E_R\times E_U)\\
&\quad=
\left[
\int_{\mathcal O_d^n}\int_{\mathcal R}
\mathbf1_{E_R}(\mathsf r)\,
\mathsf L_{\mathrm{tr}}(\mathbf o,\mathsf r)(E_Z)\,
\lambda(d\mathsf r)\,P_O^{\otimes n}(d\mathbf o)
\right]
\operatorname{Uniform}[0,1](E_U).
\end{aligned}
\]
Indeed, for fixed \(\mathbf o,\mathsf r,\mathsf u\), the rectangle indicator contributes the transcript probability multiplied by
\(\mathbf1_{E_R}(\mathsf r)\mathbf1_{E_U}(\mathsf u)\); integrating in \(\mathsf u\) gives the displayed factor. Taking \(E_U=\mathbb R\) identifies the bracket as the transcript-seed marginal. More precisely, define
\[
\mathsf L_{\mathrm{dec}}^{ZR}(P)(E_{ZR})
:=
\mathsf L_{\mathrm{dec}}(P)
\{(\mathbf z,\mathsf r,\mathsf u):
(\mathbf z,\mathsf r)\in E_{ZR}\}.
\]
Equality on measurable rectangles therefore proves, with the product identified with decision triples,
\[
\mathsf L_{\mathrm{dec}}(P)
=
\mathsf L_{\mathrm{dec}}^{ZR}(P)
\otimes\operatorname{Uniform}[0,1],
\qquad P\in\mathcal V_d.
\]
% lean: CausalSmith.Stat.LdpOptvalueUniformFrontier.sampling_decisionLaw_randomizer

\item \textit{Total variation of the decision mixtures.}
Define a common probability kernel that appends the analyst randomizer:
\[
\mathsf C_{\mathrm{coin}}((\mathsf r,\mathbf z),E)
:=
\int_{[0,1]}
\mathbf1_E(\mathbf z,\mathsf r,\mathsf u)\,d\mathsf u.
\]
The rowwise product identity from the preceding step implies
\[
\mathsf L_{\mathrm{dec}}(P)(E)
=
\int_{\mathcal Z\times\mathcal R}
\mathsf C_{\mathrm{coin}}((\mathsf r,\mathbf z),E)\,
\mathsf L_{\mathrm{dec}}^{ZR}(P)(d\mathbf z,d\mathsf r),
\qquad P\in\mathcal V_d.
\]
Integrating this equality under either supported prior and interchanging the nonnegative integrals gives
\[
\mathsf L_\ell(E)
=
\int_{\mathcal R\times\mathcal Z}
\mathsf C_{\mathrm{coin}}((\mathsf r,\mathbf z),E)\,
N_\ell(d\mathsf r,d\mathbf z).
\]

For a Borel event \(E\), put
\[
f_E(\mathsf r,\mathbf z)
:=
\mathsf C_{\mathrm{coin}}((\mathsf r,\mathbf z),E),
\qquad
0\le f_E\le1.
\]
The function is measurable. Its layer-set representation gives
\[
\begin{aligned}
|\mathsf L_0(E)-\mathsf L_1(E)|
&=
\left|
\int_0^1
\left[
N_0\{f_E\ge t_{\mathrm{lev}}\}
-
N_1\{f_E\ge t_{\mathrm{lev}}\}
\right]dt_{\mathrm{lev}}
\right|\\
&\le
\int_0^1
\left|
N_0\{f_E\ge t_{\mathrm{lev}}\}
-
N_1\{f_E\ge t_{\mathrm{lev}}\}
\right|dt_{\mathrm{lev}}\\
&\le \operatorname{TV}(N_0,N_1).
\end{aligned}
\]
Taking the supremum over Borel decision events proves
\[
\operatorname{TV}(\mathsf L_0,\mathsf L_1)
\le \operatorname{TV}(N_0,N_1)
\le\omega.
\]
% lean: CausalSmith.Stat.LdpOptvalueUniformFrontier.supported_mixture_randomizer_tv

\item \textit{Probability rows for the testing argument.}
Since a probability prior is supported on \(\mathcal V_d\), that class is nonempty. Choose
\[
P_\star\in\mathcal V_d.
\]
Complete the measurable decision laws to probability rows on the entire auxiliary parameter space by setting
\[
\mathsf C_{\mathrm{dec}}(P):=
\begin{cases}
\mathsf L_{\mathrm{dec}}(P),
&\mathsf L_{\mathrm{dec}}(P)
(\mathcal Z\times\mathcal R\times\mathbb R)=1,\\
\mathsf L_{\mathrm{dec}}(P_\star),
&\text{otherwise}.
\end{cases}
\]
The condition selecting the first branch is measurable. The resulting kernel has probability rows, and
\[
\mathsf C_{\mathrm{dec}}(P)
=\mathsf L_{\mathrm{dec}}(P)
\quad(P\in\mathcal V_d),
\qquad
\int\mathsf C_{\mathrm{dec}}(P)\,\pi_\ell(dP)
=\mathsf L_\ell.
\]
Thus the completion preserves both prior-predictive laws and their total-variation bound. The concentration and proximity hypotheses also imply
\[
\eta_0\ge0,
\qquad
\omega\ge0,
\]
because probabilities and total variation are nonnegative.
% lean: CausalSmith.Stat.LdpOptvalueUniformFrontier.separated_prior_squared_risk

\item \textit{Squared-error bound.}
Fix a Borel estimator \(T\). Define
\[
\begin{gathered}
v_{\mathrm{mid}}:=\frac{v_0+v_1}{2},
\qquad
c_{\mathrm{err}}:=\left(\frac{3\Delta}{8}\right)^2,\\
E_T:=\{\mathsf w_{\mathrm{dec}}:
v_{\mathrm{mid}}\le T(\mathsf w_{\mathrm{dec}})\},
\qquad
E_{\mathrm{err},0}:=E_T,
\qquad
E_{\mathrm{err},1}:=E_T^c,\\
E_{\mathrm{esc},\ell}
:=\{P:\Delta/8<|V(P)-v_\ell|\},
\qquad \ell\in\{0,1\},\\
\mathcal E_T(P)
:=
\int
(T(\mathsf w_{\mathrm{dec}})-V(P))^2\,
\mathsf C_{\mathrm{dec}}(P)(d\mathsf w_{\mathrm{dec}}),
\qquad
\mathcal B_{\ell,T}
:=
\int\mathcal E_T(P)\,\pi_\ell(dP).
\end{gathered}
\]
The last two integrals are nonnegative integrals with values in
\([0,\infty]\).

For each prior index, the following pointwise inequality holds over the entire auxiliary parameter space:
\[
c_{\mathrm{err}}\,
\mathsf C_{\mathrm{dec}}(P)(E_{\mathrm{err},\ell})
\le
\mathcal E_T(P)
+
c_{\mathrm{err}}\mathbf1_{E_{\mathrm{esc},\ell}}(P).
\]
To see this, first suppose \(P\in E_{\mathrm{esc},\ell}\). The left side is at most \(c_{\mathrm{err}}\), since the row is a probability measure. Otherwise,
\(|V(P)-v_\ell|\le\Delta/8\). For \(\ell=0\), on \(E_T\),
\[
T(\mathsf w_{\mathrm{dec}})-V(P)
\ge v_{\mathrm{mid}}-(v_0+\Delta/8)
=\frac{3\Delta}{8}.
\]
For \(\ell=1\), on \(E_T^c\),
\[
V(P)-T(\mathsf w_{\mathrm{dec}})
>
(v_1-\Delta/8)-v_{\mathrm{mid}}
=\frac{3\Delta}{8}.
\]
In either case the squared loss dominates
\(c_{\mathrm{err}}\mathbf1_{E_{\mathrm{err},\ell}}\), proving the pointwise inequality by integration.

Averaging under each prior gives
\[
c_{\mathrm{err}}\mathsf L_0(E_T)
\le\mathcal B_{0,T}+c_{\mathrm{err}}\pi_0(E_{\mathrm{esc},0}),
\qquad
c_{\mathrm{err}}\mathsf L_1(E_T^c)
\le\mathcal B_{1,T}+c_{\mathrm{err}}\pi_1(E_{\mathrm{esc},1}).
\]
The testing inequality follows directly from total variation:
\[
\mathsf L_0(E_T)+\mathsf L_1(E_T^c)
=
1+\mathsf L_0(E_T)-\mathsf L_1(E_T)
\ge1-\omega.
\]
If \(\max\{\mathcal B_{0,T},\mathcal B_{1,T}\}=\infty\), the Bayes-risk lower bound below follows immediately. Otherwise, both Bayes risks are finite, so summing the preceding inequalities and using the escape bounds yields
\[
\begin{aligned}
c_{\mathrm{err}}(1-\omega-2\eta_0)
&\le \mathcal B_{0,T}+\mathcal B_{1,T}\\
&\le2\max\{\mathcal B_{0,T},\mathcal B_{1,T}\}.
\end{aligned}
\]
Consequently, in both cases,
\[
\frac{9\Delta^2}{128}(1-\omega-2\eta_0)
\le\max\{\mathcal B_{0,T},\mathcal B_{1,T}\}.
\]

Each supported prior average is bounded by the supremum on the causal class:
\[
\mathcal B_{\ell,T}
=
\int_{\mathcal V_d}\mathcal E_T(P)\,\pi_\ell(dP)
\le
\sup_{P\in\mathcal V_d}
\int
(T(\mathsf w_{\mathrm{dec}})-V(P))^2\,
\mathsf L_{\mathrm{dec}}(P)(d\mathsf w_{\mathrm{dec}}).
\]
If \(1-\omega-2\eta_0\ge0\), combining these inequalities gives the stated positive-part bound. If \(1-\omega-2\eta_0<0\), its positive part is zero and the bound follows from nonnegativity of squared risk. Finally, equality of the decision laws on \(\mathcal V_d\) identifies the supremum:
\[
\begin{aligned}
&\sup_{P\in\mathcal V_d}
\int
(T(\mathsf w_{\mathrm{dec}})-V(P))^2\,
\mathsf L_{\mathrm{dec}}(P)(d\mathsf w_{\mathrm{dec}})\\
&\qquad=
\sup_{P\in\mathcal V_d}
\mathbb E_{P,Q}\bigl[(T(\mathbf Z,R,U)-V(P))^2\bigr].
\end{aligned}
\]
This proves the risk assertion with its extended-nonnegative interpretation.
% lean: CausalSmith.Stat.LdpOptvalueUniformFrontier.separated_prior_squared_risk

\item \textit{Coverage of the target neighborhoods.}
Fix an interval satisfying the stated endpoint and coverage conditions. Write
\[
I(\mathsf w_{\mathrm{dec}})
=
[\underline I(\mathsf w_{\mathrm{dec}}),
 \overline I(\mathsf w_{\mathrm{dec}})],
\qquad
\operatorname{length}I(\mathsf w_{\mathrm{dec}})
=
\overline I(\mathsf w_{\mathrm{dec}})
-\underline I(\mathsf w_{\mathrm{dec}}).
\]
For \(\ell\in\{0,1\}\), define
\[
J_\ell:=[v_\ell-\Delta/8,v_\ell+\Delta/8],
\qquad
E_{I,\ell}:=
\{\mathsf w_{\mathrm{dec}}:
\underline I(\mathsf w_{\mathrm{dec}})\le v_\ell+\Delta/8,\
v_\ell-\Delta/8\le\overline I(\mathsf w_{\mathrm{dec}})\}.
\]
These are Borel events, and ordered endpoints make them exactly the events
\(I\cap J_\ell\ne\varnothing\).

For \(P\in\mathcal V_d\setminus E_{\mathrm{esc},\ell}\),
\[
\{\,\mathsf w_{\mathrm{dec}}:
V(P)\in I(\mathsf w_{\mathrm{dec}})\,\}
\subseteq E_{I,\ell},
\]
since \(V(P)\in J_\ell\). Global coverage therefore gives
\[
0.90
\le
\mathsf L_{\mathrm{dec}}(P)(E_{I,\ell})
+\mathbf1_{E_{\mathrm{esc},\ell}}(P),
\qquad P\in\mathcal V_d.
\]
For a law in the escape event, this inequality follows from the indicator being one and the event probability being nonnegative; for a law outside it, it follows from the displayed inclusion and coverage. Integrating under the supported prior proves
\[
0.90
\le
\mathsf L_\ell(E_{I,\ell})
+\pi_\ell(E_{\mathrm{esc},\ell}),
\qquad
\mathsf L_\ell(E_{I,\ell})\ge0.90-\eta_0.
\]
% lean: CausalSmith.Stat.LdpOptvalueUniformFrontier.supported_prior_neighborhood_mass

\item \textit{Intersection and expected length.}
The total-variation bound transfers the second neighborhood event to the first mixture:
\[
\mathsf L_0(E_{I,1})
\ge\mathsf L_1(E_{I,1})-\omega
\ge0.90-\eta_0-\omega.
\]
Under the first mixture,
\[
\begin{aligned}
\mathsf L_0(E_{I,0}\cap E_{I,1})
&=
\mathsf L_0(E_{I,0})
-\mathsf L_0(E_{I,0}\setminus E_{I,1})\\
&\ge
\mathsf L_0(E_{I,0})-\mathsf L_0(E_{I,1}^c)\\
&=
\mathsf L_0(E_{I,0})+\mathsf L_0(E_{I,1})-1\\
&\ge0.80-2\eta_0-\omega.
\end{aligned}
\]
Combining this with nonnegativity of probability gives
\[
\mathsf L_0(E_{I,0}\cap E_{I,1})
\ge(0.80-2\eta_0-\omega)_+.
\]

On the intersection event the endpoints satisfy
\[
\underline I(\mathsf w_{\mathrm{dec}})\le v_0+\Delta/8,
\qquad
\overline I(\mathsf w_{\mathrm{dec}})\ge v_1-\Delta/8.
\]
Hence
\[
\operatorname{length}I(\mathsf w_{\mathrm{dec}})
\ge
(v_1-\Delta/8)-(v_0+\Delta/8)
=\frac{3\Delta}{4}.
\]
Outside that event, ordered endpoints give nonnegative length. Thus the pointwise indicator bound and its integrated form are
\[
\operatorname{length}I(\mathsf w_{\mathrm{dec}})
\ge
\frac{3\Delta}{4}\,
\mathbf1_{E_{I,0}\cap E_{I,1}}(\mathsf w_{\mathrm{dec}}),
\]
\[
\int\operatorname{length}I\,d\mathsf L_0
\ge
\frac{3\Delta}{4}\,
\mathsf L_0(E_{I,0}\cap E_{I,1})
\ge
\frac{3\Delta}{4}(0.80-2\eta_0-\omega)_+.
\]
Finally, nonnegative integration through the mixture and causal support yield
\[
\begin{aligned}
\int\operatorname{length}I\,d\mathsf L_0
&=
\int_{\mathcal V_d}
\left[
\int\operatorname{length}I\,d\mathsf L_{\mathrm{dec}}(P)
\right]\pi_0(dP)\\
&\le
\sup_{P\in\mathcal V_d}
\int\operatorname{length}I\,d\mathsf L_{\mathrm{dec}}(P)\\
&=
\sup_{P\in\mathcal V_d}
\mathbb E_{P,Q}\operatorname{length}I(\mathbf Z,R,U).
\end{aligned}
\]
The same equality of decision laws transfers the global coverage condition to every completed row on \(\mathcal V_d\), so this establishes the interval assertion.
% lean: CausalSmith.Stat.LdpOptvalueUniformFrontier.separated_prior_interval_length
\end{enumerate}
\end{proof}
